\section{Einführung und Ziele}

\subsection{Aufgabenstellung}

Die Aufgabenstellung bestand darin, einen Architekturstil anhand einer
Beispielanwendung umzusetzen.
Dabei fiel die Wahl auf die Microservice-Architektur.

Als Beispielanwendung dient eine E-Commerce-Plattform,
deren Funktionen sich gut in fachlich getrennte Bereiche aufteilen lassen.
Der vollständige Quellcode ist im zugehörigen GitHub-Repository verfügbar:

\begin{center}
    \url{https://github.com/noah-preyer/SS26_Software_Architektur_Ebay2.git}
\end{center}

Die Umsetzung konzentriert sich auf die folgenden Kernfunktionen:

\subsection{Kernfunktionen}

\begin{itemize}
\item Registrierung und Anmeldung mittels JWT
\item Erstellen, Bearbeiten, Löschen und Suchen von Produkten
\item beispielhafte Kaufabwicklung und E-Mail-Benachrichtigungen
\end{itemize}


\subsection{Qualitätsziele}

\begin{table}[ht]
    \centering
    {
        \renewcommand{\arraystretch}{1.5}
        \rowcolors{2}{tableRowColor}{white}

        \begin{tabularx}{\linewidth}{
                >{\bfseries\raggedright\arraybackslash}p{0.20\linewidth}
                >{\raggedright\arraybackslash}X
        }
        \toprule
        Qualitätsziel & Beschreibung \\
        \midrule

        Modularität &
        Die Services sollen fachlich klar voneinander abgegrenzt sein. \\

        Sicherheit &
        Geschützte Service-Funktionen sollen nur für berechtigte Benutzer zugänglich sein. \\

        Erweiterbarkeit &
        Neue Services sollen mit möglichst wenigen Änderungen am bestehenden System ergänzt werden können. \\

        Skalierbarkeit &
        Ein Service soll bei Bedarf unabhängig vom restlichen System durch mehrere Replikate skaliert werden können. \\

        Fehlertoleranz &
        Der Ausfall eines Services soll andere Bereiche möglichst nicht
        beeinträchtigen. \\

        \bottomrule
        \end{tabularx}
    }

    \caption{Qualitätsziele der Architektur}
    \label{tab:qualitaetsziele}
\end{table}